\section{Threats to Validity}
\label{sec:threats}

\paragraph{Construct validity.}
MIC and HC50 are assay-dependent measurements.
Consensus labels cannot remove all variation from strain, medium, peptide purity,
exposure time, or readout protocol.
The GRAMPA experiments use $\log_2$ MIC values, whereas QMAP evaluates
$\log_{10}$ consensus values through its official metric interface.
These targets are appropriate for regression, but they are not direct clinical
efficacy or safety endpoints.
HC50 also captures hemolysis, not the full toxicity profile needed for therapy.

\paragraph{Internal validity.}
For QMAP we use the official five predefined splits, split-specific
\texttt{get\_train\_mask} filtering, and \texttt{compute\_metrics(..., log=True)}
to avoid silent split drift.
The main JEPA QMAP claims are averaged over three training seeds, reducing the
risk that a single favorable initialization drives the result.
Still, the conditional model pools multiple bacterial targets and HC50 under
shared property keys.
Its gains should therefore be read as evidence for useful shared conditioning,
not as proof that one universal head is best for every endpoint.

\paragraph{External validity.}
The current JEPA encoder supports peptides up to 50 amino acids after reserving
special tokens, so longer QMAP sequences are truncated during evaluation.
QMAP includes noncanonical amino acids and SMILES representations for some
peptides; the current sequence encoder treats the benchmark as a sequence-only
task and does not yet model full chemical structure, terminal modifications, or
intrachain bonds explicitly.
This makes the reported QMAP performance conservative for sequence-level
comparison, but it also limits claims about chemically modified peptides.

\paragraph{Conclusion validity.}
Leaderboard comparisons are restricted to the methods and numbers reported in
the QMAP README plus our reproduced JEPA runs.
We report seed variability, but we do not yet provide paired statistical tests
over split-level predictions against every prior method because their raw
predictions are not available in this workspace.
Finally, generated peptides are evaluated with computational oracles and
physicochemical checks only.
Wet-lab synthesis and antimicrobial assays are still required before making
biological efficacy claims.

\paragraph{Dry-lab scope.}
All claims in this manuscript are computational. The MIC-conditioned generator is
therefore evaluated as an oracle-mediated steering mechanism, not as evidence
that generated sequences are active, selective, safe, stable, or synthesizable.
Future wet-lab validation would change the evidentiary status of candidate
claims, but the current paper should be read as a benchmarked dry-lab study.
